% Two ways of organizing shared memory among several processes.
% Usage: % Two ways of organizing shared memory among several processes.
% Usage: % Two ways of organizing shared memory among several processes.
% Usage: % Two ways of organizing shared memory among several processes.
% Usage: \input{figures/l09-segments}   (width ~118 mm)
\begin{tikzpicture}[x=1mm, y=1mm, font=\footnotesize\sffamily,
  blk/.style={draw, minimum height=8mm, inner sep=0pt, font=\scriptsize\sffamily},
  used/.style={blk, fill=ln-accent!20},
  freeb/.style={blk, fill=white},
  box/.style={draw, rounded corners=2pt, fill=ln-box, align=center,
              inner sep=1.5pt, font=\scriptsize\sffamily},
  cell/.style={draw, fill=white, minimum width=5mm, minimum height=5mm,
               inner sep=0pt, font=\scriptsize\sffamily},
  lab/.style={font=\scriptsize\sffamily, text=ln-muted, align=center},
  >={Stealth[length=1.6mm]}]
  % (a) one large segment managed by an allocator
  \node[used,  minimum width=12mm, anchor=west] (a1) at (0,0) {P1/P2};
  \node[freeb, minimum width=8mm,  anchor=west] (a2) at (a1.east) {};
  \node[used,  minimum width=12mm, anchor=west] (a3) at (a2.east) {P1/P3};
  \node[used,  minimum width=12mm, anchor=west] (a4) at (a3.east) {P2/P3};
  \node[freeb, minimum width=8mm,  anchor=west] (a5) at (a4.east) {};
  \node[box, text width=24mm] (alloc) at (26,18)
    {\iflangja{アロケータ\\（確保・解放）}{allocator\\(allocate / free)}};
  \foreach \t in {a1,a3,a4} \draw[->] (alloc.south) -- (\t.north);
  \node[lab] at (26,-9) {\iflangja{(a) 1つの大きなセグメント}{(a) one large segment}};
  % (b) pool of small segments
  \foreach \i/\s in {1/used,2/freeb,3/used,4/freeb} {
    \node[\s, minimum width=10mm] (s\i) at (58+13*\i,0) {seg \i};
  }
  \node[lab, anchor=east] at (86.5,18) {\iflangja{空き ID}{free ids}};
  \node[cell] (q2) at (90,18) {2};
  \node[cell] (q4) at (95,18) {4};
  \draw[->, dashed] (q2.south) -- (s2.north);
  \draw[->, dashed] (q4.south) -- (s4.north);
  \node[lab] at (90.5,-9) {\iflangja{(b) 小さなセグメントのプール}{(b) pool of small segments}};
\end{tikzpicture}
   (width ~118 mm)
\begin{tikzpicture}[x=1mm, y=1mm, font=\footnotesize\sffamily,
  blk/.style={draw, minimum height=8mm, inner sep=0pt, font=\scriptsize\sffamily},
  used/.style={blk, fill=ln-accent!20},
  freeb/.style={blk, fill=white},
  box/.style={draw, rounded corners=2pt, fill=ln-box, align=center,
              inner sep=1.5pt, font=\scriptsize\sffamily},
  cell/.style={draw, fill=white, minimum width=5mm, minimum height=5mm,
               inner sep=0pt, font=\scriptsize\sffamily},
  lab/.style={font=\scriptsize\sffamily, text=ln-muted, align=center},
  >={Stealth[length=1.6mm]}]
  % (a) one large segment managed by an allocator
  \node[used,  minimum width=12mm, anchor=west] (a1) at (0,0) {P1/P2};
  \node[freeb, minimum width=8mm,  anchor=west] (a2) at (a1.east) {};
  \node[used,  minimum width=12mm, anchor=west] (a3) at (a2.east) {P1/P3};
  \node[used,  minimum width=12mm, anchor=west] (a4) at (a3.east) {P2/P3};
  \node[freeb, minimum width=8mm,  anchor=west] (a5) at (a4.east) {};
  \node[box, text width=24mm] (alloc) at (26,18)
    {\iflangja{アロケータ\\（確保・解放）}{allocator\\(allocate / free)}};
  \foreach \t in {a1,a3,a4} \draw[->] (alloc.south) -- (\t.north);
  \node[lab] at (26,-9) {\iflangja{(a) 1つの大きなセグメント}{(a) one large segment}};
  % (b) pool of small segments
  \foreach \i/\s in {1/used,2/freeb,3/used,4/freeb} {
    \node[\s, minimum width=10mm] (s\i) at (58+13*\i,0) {seg \i};
  }
  \node[lab, anchor=east] at (86.5,18) {\iflangja{空き ID}{free ids}};
  \node[cell] (q2) at (90,18) {2};
  \node[cell] (q4) at (95,18) {4};
  \draw[->, dashed] (q2.south) -- (s2.north);
  \draw[->, dashed] (q4.south) -- (s4.north);
  \node[lab] at (90.5,-9) {\iflangja{(b) 小さなセグメントのプール}{(b) pool of small segments}};
\end{tikzpicture}
   (width ~118 mm)
\begin{tikzpicture}[x=1mm, y=1mm, font=\footnotesize\sffamily,
  blk/.style={draw, minimum height=8mm, inner sep=0pt, font=\scriptsize\sffamily},
  used/.style={blk, fill=ln-accent!20},
  freeb/.style={blk, fill=white},
  box/.style={draw, rounded corners=2pt, fill=ln-box, align=center,
              inner sep=1.5pt, font=\scriptsize\sffamily},
  cell/.style={draw, fill=white, minimum width=5mm, minimum height=5mm,
               inner sep=0pt, font=\scriptsize\sffamily},
  lab/.style={font=\scriptsize\sffamily, text=ln-muted, align=center},
  >={Stealth[length=1.6mm]}]
  % (a) one large segment managed by an allocator
  \node[used,  minimum width=12mm, anchor=west] (a1) at (0,0) {P1/P2};
  \node[freeb, minimum width=8mm,  anchor=west] (a2) at (a1.east) {};
  \node[used,  minimum width=12mm, anchor=west] (a3) at (a2.east) {P1/P3};
  \node[used,  minimum width=12mm, anchor=west] (a4) at (a3.east) {P2/P3};
  \node[freeb, minimum width=8mm,  anchor=west] (a5) at (a4.east) {};
  \node[box, text width=24mm] (alloc) at (26,18)
    {\iflangja{アロケータ\\（確保・解放）}{allocator\\(allocate / free)}};
  \foreach \t in {a1,a3,a4} \draw[->] (alloc.south) -- (\t.north);
  \node[lab] at (26,-9) {\iflangja{(a) 1つの大きなセグメント}{(a) one large segment}};
  % (b) pool of small segments
  \foreach \i/\s in {1/used,2/freeb,3/used,4/freeb} {
    \node[\s, minimum width=10mm] (s\i) at (58+13*\i,0) {seg \i};
  }
  \node[lab, anchor=east] at (86.5,18) {\iflangja{空き ID}{free ids}};
  \node[cell] (q2) at (90,18) {2};
  \node[cell] (q4) at (95,18) {4};
  \draw[->, dashed] (q2.south) -- (s2.north);
  \draw[->, dashed] (q4.south) -- (s4.north);
  \node[lab] at (90.5,-9) {\iflangja{(b) 小さなセグメントのプール}{(b) pool of small segments}};
\end{tikzpicture}
   (width ~118 mm)
\begin{tikzpicture}[x=1mm, y=1mm, font=\footnotesize\sffamily,
  blk/.style={draw, minimum height=8mm, inner sep=0pt, font=\scriptsize\sffamily},
  used/.style={blk, fill=ln-accent!20},
  freeb/.style={blk, fill=white},
  box/.style={draw, rounded corners=2pt, fill=ln-box, align=center,
              inner sep=1.5pt, font=\scriptsize\sffamily},
  cell/.style={draw, fill=white, minimum width=5mm, minimum height=5mm,
               inner sep=0pt, font=\scriptsize\sffamily},
  lab/.style={font=\scriptsize\sffamily, text=ln-muted, align=center},
  >={Stealth[length=1.6mm]}]
  % (a) one large segment managed by an allocator
  \node[used,  minimum width=12mm, anchor=west] (a1) at (0,0) {P1/P2};
  \node[freeb, minimum width=8mm,  anchor=west] (a2) at (a1.east) {};
  \node[used,  minimum width=12mm, anchor=west] (a3) at (a2.east) {P1/P3};
  \node[used,  minimum width=12mm, anchor=west] (a4) at (a3.east) {P2/P3};
  \node[freeb, minimum width=8mm,  anchor=west] (a5) at (a4.east) {};
  \node[box, text width=24mm] (alloc) at (26,18)
    {\iflangja{アロケータ\\（確保・解放）}{allocator\\(allocate / free)}};
  \foreach \t in {a1,a3,a4} \draw[->] (alloc.south) -- (\t.north);
  \node[lab] at (26,-9) {\iflangja{(a) 1つの大きなセグメント}{(a) one large segment}};
  % (b) pool of small segments
  \foreach \i/\s in {1/used,2/freeb,3/used,4/freeb} {
    \node[\s, minimum width=10mm] (s\i) at (58+13*\i,0) {seg \i};
  }
  \node[lab, anchor=east] at (86.5,18) {\iflangja{空き ID}{free ids}};
  \node[cell] (q2) at (90,18) {2};
  \node[cell] (q4) at (95,18) {4};
  \draw[->, dashed] (q2.south) -- (s2.north);
  \draw[->, dashed] (q4.south) -- (s4.north);
  \node[lab] at (90.5,-9) {\iflangja{(b) 小さなセグメントのプール}{(b) pool of small segments}};
\end{tikzpicture}
